LLM coding agents read logs, configuration files and tickets through tools, and whatever a tool returns enters the model's context. Agents are increasingly given local redaction tools through the Model Context Protocol, so we ask what an agent learns through such a tool. In a pre-registered study of an open-source on-device guardrail, a differential audit planted values, called every tool, and searched every reply in five encodings against a masked control. It found that a reply documented as content-free returned case-folded person names through detector metadata: with a name-list tier active, in 701 of 2,000 generated agent tool outputs and 24 of 129 real WNUT-17 documents. An allow-list of value-free metadata keys removed every hit. A regex-gated INT8 DistilBERT tier raised person-name character recall on real user-generated text (WNUT-17) by 0.482 over the pattern tiers. The gate, which sends 30\% of the text to the model, cost 0.225 of recall against running the model on every line. Batched inference broke streaming equivalence and left up to 12 names in clear that a whole-document pass redacted. Merged windows exceeded the tier's declared 256-character context. Measuring per-call cost exposed a Nagle stall on every kept-alive REST request; fixing it cut a 256-byte call from 44.00 to 0.75\,ms.
